\begin{center}
    \textbf{\uppercase{Resumo}}
\end{center}
\vspace*{1.5cm}


Esta pesquisa analisa a dinâmica dos estados brasileiros após um choque de incerteza agregada, utilizando duas medidas de incerteza doméstica e uma de incerteza global. Para isto foi utilizado um modelo \gls{SVAR} conforme sugerido por \citeonline{baker2016measuring} e \citeonline{barboza2018efeitos}, para obter a \gls{FIR} de cada estado. Posteriormente, foi estimado um modelo de regressão linear múltipla com corte transversal para relacionar as respostas dos estados e suas respectivas estruturas econômicas. Os resultados indicam que os estados reagem de forma heterogênea independente de medida de incerteza ou de atividade econômica, porém medidas de incerteza domésticas relacionadas a expectativas de mercado possuem maior impacto negativo na economia, além do mais, a produção industrial tende a ser mais afetada quando comparada a \textit{proxy} do \gls{PIB}, o \gls{IBCR}. Por fim, estados com com maior participação de setores como agropecuária, indústria extrativa e de transformação, serviços financeiros, serviços públicos e maior grau de abertura relativos ao \gls{PIB} do estado, são menos impactados, no entanto, onde há maior participação de ramos de construção e alto grau de orçamento público, são mais fortemente afetados.

\vspace*{1.5cm}
\noindent \textbf{Palavras-chave:}  macroeconomia; choque de incerteza; atividade econômica; séries temporais, SVAR, ciclos de negócios.

\thispagestyle{empty}